\documentclass[leqno]{article}
% Shared QBT/Soln pack preamble for resources.danieljsmith.org.
\usepackage{graphicx}
\usepackage{dsfont}
\usepackage{mathtools}
\usepackage{amssymb}
\usepackage{amsmath}
\usepackage{mleftright}
\usepackage{amsthm}
\usepackage{float}
\usepackage{ragged2e}
\usepackage{tensor}
\usepackage{fancyhdr}
\usepackage{empheq}
\usepackage[most]{tcolorbox}
\usepackage{enumitem}
\usepackage{dirtytalk}
\usepackage{marginnote}
\usepackage{microtype}
\usepackage[
  a4paper,
  left=0.8in,
  right=1in,
  top=1in,
  bottom=0.5in,
  headheight=14pt,
  headsep=0.4in
]{geometry}
\setlength{\parindent}{0cm}

% TikZ / PGFPlots
\usepackage{pgfplots}
\pgfplotsset{compat=1.18}
\usepgfplotslibrary{fillbetween, polar}
\usetikzlibrary{calc, positioning}
\usetikzlibrary{angles, quotes}
\usetikzlibrary{arrows.meta}
\usetikzlibrary{decorations.pathreplacing, decorations.markings}
\usetikzlibrary{patterns}
\usetikzlibrary{intersections}

% Links and contents. The contents follow the TMUA papers: compact entries,
% no dot leaders, and each question's mark total in grey after its label.
\usepackage{tocloft}
\tocloftpagestyle{fancy}
\setlength{\cftbeforesubsecskip}{2pt}
\setlength{\cftsubsecindent}{1.5em}
\setlength{\cftsubsecnumwidth}{0pt}
\renewcommand{\cftsubsecfont}{\small\color{black}}
\renewcommand{\cftsubsecpagefont}{\small\color{black}}
\renewcommand{\cftsubsecleader}{\hfill}
\usepackage[colorlinks=true,linkcolor=black,citecolor=magenta,urlcolor=black]{hyperref}%
\urlstyle{same}

% Pack macros
% \marks also adds its value to the current question's total in the .aux
% file, so the contents show each total from the second pass. The write
% happens at shipout, so a box typeset more than once counts once. The
% pack's mark total for the contents-page title is summed the same way.
\makeatletter
\newcount\djs@totalmarks
\newcommand{\djsAddMarks}[2]{%
  \global\advance\djs@totalmarks#2\relax
  \@ifundefined{djs@marks@#1}%
    {\expandafter\xdef\csname djs@marks@#1\endcsname{#2}}%
    {\expandafter\xdef\csname djs@marks@#1\endcsname
      {\the\numexpr\csname djs@marks@#1\endcsname+#2\relax}}}
\RenewDocumentCommand{\marks}{o m}{%
  \protected@write\@auxout{}{\string\djsAddMarks{\number\value{enumi}}{#2}}%
  \IfValueTF{#1}%
    {\marginnote{\raggedright\textbf{[#2]}}[#1]}%
    {\marginnote{\raggedright\textbf{[#2]}}}%
}
\makeatother

\newcommand{\vspacerule}[2]{%
  \par\vspace*{#1}%
  \noindent\hrulefill\par
  \vspace*{#2}%
}

\definecolor{scarlet}{RGB}{205, 40, 75}
% Site colours for the contents-page title: the accent, the rule colour, a
% stronger rule colour for the path separator, and the muted text colour.
\definecolor{djsBurgundy}{RGB}{118, 59, 67}
\definecolor{djsHairline}{RGB}{216, 213, 205}
\definecolor{djsSlash}{RGB}{170, 166, 158}
\definecolor{djsMuted}{RGB}{96, 101, 106}

% Worked-solution box, adapted from the TMUA worked-solutions box. A white breakable box with a left scarlet rule and no surrounding
% frame or tint. A part that breaks fills the rest of its page, so the rule
% runs to the bottom margin; the unbroken or last part ends in a short
% horizontal foot rule so a solution visibly stops. Unlike the TMUA box there
% is no answer label. The box does not grow to the left, so the rule sits at
% the text edge; growing to the right keeps the usable line width close to the
% old tcolorbox Solution box.
% Spacing inside the box: paragraphs are 6pt apart, so source blank lines show
% (solutions put one after each display and none between a lead-in and its
% display). Every display gets 6pt above and 10pt below, whatever the length
% of its lead-in line; align rows are spread by an 8pt \jot; solution parts
% are 14pt apart and sub-parts 8pt; \left and \right space like ordinary
% brackets (mleftright); no single line is stranded at a page break; and
% neither the Solution label nor the line before a list ends a page. No
% skip can shrink, so full pages are not squeezed.
\newtcolorbox{djsSolution}{%
  enhanced, breakable, frame hidden, sharp corners, colback=white,
  height fixed for=first and middle, height fill,
  extras unbroken and last={natural height},
  borderline west={2.2pt}{0pt}{scarlet},
  left=11pt, right=0pt, top=5pt, bottom=13pt,
  before skip=13pt, after skip=4pt,
  grow to right by=10mm,
  overlay unbroken and last={%
    \fill[scarlet] (frame.south west)
      rectangle ([xshift=3.5cm, yshift=2.2pt]frame.south west);},
  before upper={{\color{scarlet}\bfseries Solution}\par\nobreak\vspace{4pt}%
    \setlength{\parskip}{6pt plus 2pt}%
    \setlength{\abovedisplayskip}{6pt plus 2pt}%
    \setlength{\abovedisplayshortskip}{6pt plus 2pt}%
    \setlength{\belowdisplayskip}{10pt plus 2pt}%
    \setlength{\belowdisplayshortskip}{10pt plus 2pt}%
    \setlength{\jot}{8pt}%
    \setlist{parsep=6pt plus 2pt, topsep=6pt plus 2pt, beginpenalty=10000}%
    \setlist[enumerate,1]{itemsep=14pt plus 2pt}%
    \setlist[enumerate,2]{itemsep=8pt plus 2pt}%
    \mleftright
    \clubpenalty=10000 \widowpenalty=10000 \displaywidowpenalty=10000\relax}}

% Optional smaller figure inside a solution box. \scalebox shrinks node text
% with the drawing. Not applied to existing figures.
\newcommand{\SolnFigureScale}{0.85}
\newsavebox{\djsSolnFigureBox}
\newenvironment{djsSolutionFigure}
  {\par\vspace{4pt}\begin{center}\begin{lrbox}{\djsSolnFigureBox}}
  {\end{lrbox}\scalebox{\SolnFigureScale}{\usebox{\djsSolnFigureBox}}%
   \end{center}\vspace{2pt}\par}

\newcommand{\cxl}[2]{%
  \tikz[baseline=(X.base)]{%
    \node[inner sep=0pt, outer sep=0pt, text=#1] (X) {$#2$};
    \draw[#1, line width=0.4pt] (X.south west)--(X.north east);
  }%
}

% Page-style hook run at the contents page and each question's opening page.
% A no-op for QBT packs; \djsSolnHeader points it at the fancy header.
\newcommand{\djsQuestionPageStyle}{}

% A contents entry: "Question N" in a box wide enough for two digits, so the
% mark totals line up, then the total in grey. Bookmarks keep "Question N".
\makeatletter
\newlength{\djsTocLabelWidth}
\newcommand{\djsTocMarks}[1]{%
  \@ifundefined{djs@marks@#1}{}{%
    {\footnotesize\itshape\color{black!45}%
     \csname djs@marks@#1\endcsname~mark%
     \ifnum\csname djs@marks@#1\endcsname=1 \else s\fi}}}
\newcommand{\djsTocEntry}[1]{%
  \texorpdfstring
    {\settowidth{\djsTocLabelWidth}{Question 00}%
     \makebox[\dimexpr\djsTocLabelWidth+0.9em\relax][l]{Question #1}%
     \djsTocMarks{#1}}%
    {Question #1}}
\makeatother

% The question count for the contents-page title is written to the .aux at
% the end of the document, so it also appears from the second pass.
\newcounter{djsQuestions}
\newcommand{\questionitem}{%
  \djsQuestionPageStyle
  \item\phantomsection
  \stepcounter{djsQuestions}%
  \addcontentsline{toc}{subsection}{\protect\djsTocEntry{\number\value{enumi}}}%
}
\makeatletter
\newcommand{\djsSetQuestionCount}[1]{\gdef\djs@questioncount{#1}}
\AtEndDocument{\immediate\write\@auxout
  {\string\djsSetQuestionCount{\number\value{djsQuestions}}}}
\makeatother

\renewcommand{\contentsname}{Questions}

% \djsPackHeader[<link content tag>]{<left header>}; the tag says which
% document a click on the header link came from.
\newcommand{\djsPackHeader}[2][questions]{%
  \pagestyle{fancy}%
  \fancyhead{}%
  \fancyfoot{}%
  \renewcommand{\headrulewidth}{0.9pt}%
  \lhead{#2}%
  \chead{}%
  \rhead{\href{https://danieljsmith.org/?utm_source=pdf&utm_campaign=further-maths&utm_content=#1}{danieljsmith.org}}%
}

% The header's topic is also the contents-page title, under which a label
% says which booklet this is.
\newcommand{\djsPackTitle}{}
\newcommand{\djsBookletLabel}{}
\newcommand{\djsQbtHeader}[1]{%
  \djsPackHeader{#1}%
  \renewcommand{\djsPackTitle}{#1}%
  \renewcommand{\djsBookletLabel}{Questions by topic}%
}

% The colophon at the foot of page one, independent of the page style: the
% contents-page title's accent rule (the same \djs@accentrule, so the two
% always match), then the date on the left and a link to the site on the
% right, all inside the bottom margin. Every later page carries its number
% over the page count at the bottom right, under the header's link. The count
% is the kernel's \PreviousTotalPages, which every run already records in the
% .aux file. Sizes are fixed, so these footers match the TMUA papers'. The
% document supplies the date; recompiling never changes it.
\makeatletter
\newcommand{\djsLastUpdated}[1]{%
  \AtBeginDocument{\AddToHookNext{shipout/foreground}{\djs@colophon{#1}}}}
\newcommand{\djs@footR}[1]{%
  \put(\dimexpr\paperwidth-1in\relax,-\dimexpr\paperheight-0.32in\relax){%
    \makebox[0pt][r]{#1}}}
\newcommand{\djs@colophon}[1]{%
  \put(0.8in,-\dimexpr\paperheight-0.47in\relax){%
    \makebox[0pt][l]{\djs@accentrule{\textwidth}}}%
  \put(0.8in,-\dimexpr\paperheight-0.32in\relax){\makebox[0pt][l]{%
    \fontsize{9}{11}\selectfont Last updated #1}}%
  \djs@footR{\fontsize{7}{8.5}\selectfont\color{djsMuted}%
    Latest version at \href{https://resources.danieljsmith.org/}%
      {\color{black}resources.danieljsmith.org}}}
\AddToHook{shipout/foreground}{%
  \ifnum\value{page}>1
    \djs@footR{{\fontsize{9}{11}\selectfont\thepage}%
      {\fontsize{7}{8.5}\selectfont\color{djsMuted}%
        \hspace{0.3em}/\hspace{0.3em}\PreviousTotalPages}}%
  \fi}
\makeatother

% Solutions packs show the header only on the contents page and each
% question's opening page, so it never appears part-way through a solution.
% \thispagestyle marks the next page shipped out, which is the question's
% page because every solution question starts after a \newpage.
\newcommand{\djsSolnHeader}[1]{%
  \djsPackHeader[solutions]{\textit{Solutions} - #1}%
  \pagestyle{empty}%
  \renewcommand{\djsQuestionPageStyle}{\thispagestyle{fancy}}%
  \renewcommand{\djsPackTitle}{#1}%
  \renewcommand{\djsBookletLabel}{Worked solutions}%
}

% \djsSitePath{<strand>}{<topic>}, called once after the header, names where
% the pack sits on the site; the contents-page title prints it above the pack
% title, e.g. CORE PURE / FURTHER CALCULUS. Without it the line is left out.
\makeatletter
\newif\ifdjs@sitepath
\newcommand{\djs@strand}{}
\newcommand{\djs@section}{}
\newcommand{\djsSitePath}[2]{%
  \djs@sitepathtrue
  \renewcommand{\djs@strand}{#1}%
  \renewcommand{\djs@section}{#2}%
}
\makeatother

% \djsBoards{\djsBoard{<label>}{<code>}...}, called once after \djsSitePath,
% lists the pack's exam boards (every board in its strand for a pack for all
% boards); the contents-page title prints them under the booklet label, after
% a burgundy EXAM BOARDS key. A board's label and code never break across
% lines: a line break falls only between boards, where it replaces the
% separating dot.
\makeatletter
\newif\ifdjs@boards
\newif\ifdjs@firstboard
\newcommand{\djs@boardlist}{}
\newcommand{\djsBoards}[1]{%
  \djs@boardstrue
  \renewcommand{\djs@boardlist}{#1}%
}
\newcommand{\djsBoard}[2]{%
  \ifdjs@firstboard\global\djs@firstboardfalse\else
    \discretionary{}{}{\kern0.45em\textperiodcentered\kern0.45em}%
  \fi
  \mbox{#1 {\color{djsMuted}#2}}}
\newlength{\djs@boardkeywd}
\newcommand{\djs@boardkey}{{\color{djsBurgundy}\scriptsize\bfseries\djs@caps{Exam boards}}}
\newcommand{\djs@boardsline}{%
  \ifdjs@boards
    \vspace{16pt}%
    \settowidth{\djs@boardkeywd}{\djs@boardkey}%
    \noindent\makebox[\dimexpr\djs@boardkeywd+1.1em\relax][l]{\djs@boardkey}%
    \parbox[t]{\dimexpr\linewidth-\djs@boardkeywd-1.1em\relax}{%
      \raggedright\small\global\djs@firstboardtrue\djs@boardlist}\par
  \fi}
\makeatother

% Further Mechanics packs call \djsMechanicsPack in their preamble, so that
% \djsFrontMatter prints the default value of g above the contents. The box
% keeps its own space above it (before skip is the whole gap below the title,
% not an addition to it). Without the box, \djsGravityNote is just extra space
% before the contents heading, bringing the gap above it close to the gap
% between the box and the heading on Further Mechanics pages.
\newcommand{\djsGravityNote}{\vspace{14pt}}
\newcommand{\djsMechanicsPack}{%
  \renewcommand{\djsGravityNote}{%
    \begin{tcolorbox}[sharp corners, colback=white, colframe=black,
      boxrule=1pt, left=8pt, right=8pt, top=6pt, bottom=6pt,
      before skip=20pt]
    {\bfseries Instructions}\par\vspace{4pt}
    Unless a question states otherwise, take $g = 9.8\,\mathrm{m\,s^{-2}}$
    wherever a numerical value of $g$ is required. Solutions also give final
    answers for $g = 10\,\mathrm{m\,s^{-2}}$ where these differ.
    \end{tcolorbox}%
    \vspace{16pt}}%
}

% Contents-page title: the site path in spaced burgundy capitals, the pack
% title, a hairline led by a short scarlet rule (the solution box's foot rule),
% then the booklet label with the question and mark counts, and the exam
% boards line when the pack has one. The counts come from the .aux, so they
% appear from the second pass.
\makeatletter
\newsavebox{\djs@eyebrowbox}
\newcommand{\djs@caps}[1]{\textls[180]{\MakeUppercase{#1}}}
\newcommand{\djs@eyebrow}{%
  \color{djsBurgundy}\footnotesize\bfseries
  \djs@caps{\djs@strand}%
  {\color{djsSlash}\hspace{0.55em}/\hspace{0.55em}}%
  \djs@caps{\djs@section}}
\newcommand{\djs@dot}{\hspace{0.45em}\textperiodcentered\hspace{0.45em}}
\newcommand{\djs@counts}{%
  \@ifundefined{djs@questioncount}{}{%
    \djs@questioncount~question\ifnum\djs@questioncount=1 \else s\fi
    \ifnum\djs@totalmarks>0
      \djs@dot\the\djs@totalmarks~mark\ifnum\djs@totalmarks=1 \else s\fi
    \fi}}
% The accent rule: a hairline of the given width, led by a short scarlet rule
% (the solution box's foot rule). The contents-page title and the colophon
% both draw it, so their weights are always the same.
\newcommand{\djs@accentrule}[1]{%
  \makebox[0pt][l]{\color{djsHairline}\rule[0.9pt]{#1}{0.6pt}}%
  {\color{scarlet}\rule{3.5cm}{2.2pt}}}
\newcommand{\djsContentsTitle}{%
  \ifdjs@sitepath
    \sbox\djs@eyebrowbox{\djs@eyebrow}%
    \ifdim\wd\djs@eyebrowbox>\linewidth
      \@latex@warning@no@line{The \string\djsSitePath\space line is wider
        than the text, so it wraps. Shorten its topic}%
    \fi
    {\raggedright\djs@eyebrow\par}%
    \vspace{7pt}%
  \fi
  {\raggedright\fontsize{24.88}{28}\selectfont\djsPackTitle\par}%
  \vspace{9pt}%
  \noindent\djs@accentrule{\linewidth}\par
  \vspace{5pt}%
  {\small\color{djsMuted}\djsBookletLabel\hfill\djs@counts\par}%
  \djs@boardsline
  \vspace{8pt}%
}
\makeatother

\newcommand{\djsFrontMatter}{%
  \djsQuestionPageStyle
  \vspace*{6pt}%
  \djsContentsTitle
  \djsGravityNote
  \tableofcontents
  \newpage
}


\djsSolnHeader{Simple Harmonic Motion}
\djsSitePath{Further Mechanics}{Circular Motion \& Oscillations}
\djsBoards{\djsBoard{Edexcel FM2}{9FM0/4C}\djsBoard{OCR MEI Mech Major}{Y421}}
\djsMechanicsPack
\djsLastUpdated{6 November 2026}

\begin{document}
\djsFrontMatter
\begin{enumerate}[
  label=\textbf{\arabic*.},
  leftmargin=*,
  topsep=0pt,
  partopsep=0pt,
  parsep=0pt
]

% ---- Question 1 ----
\questionitem
% [Source: _QBT___Solns__Simple_Harmonic_Motion, Q1]

A particle $P$ is moving along the $x$-axis. At time $t$ seconds, the displacement of $P$ from the origin $O$ is $x$ metres, where
\[
x = 3\cos\left(\tfrac{1}{4}\pi t\right) + 3\sqrt{3}\sin\left(\tfrac{1}{4}\pi t\right)
\]
\\
\begin{enumerate}[label=\textbf{(\alph*)}]
\item Prove that $P$ is moving with simple harmonic motion. \marks{3} \\
\item Find the period of the motion. \marks{2} \\
\item Find the amplitude of the motion. \marks{2} \\
\item Find, in terms of $\pi$, the maximum speed of $P$. \marks{2} \\
\item Find the total time during each complete oscillation for which the speed of $P$ is greater than $\frac{3\sqrt{3}}{4}\pi\,\mathrm{m\,s^{-1}}$. \marks{6} \\
\end{enumerate}

\vspace*{10pt}

\begin{djsSolution}
\begin{enumerate}[label=\textbf{(\alph*)}, start=1]
\item Differentiating with respect to \(t\):
\begin{align*}
\dot{x}&=-\tfrac{3}{4}\pi\sin\left(\tfrac{1}{4}\pi t\right)+\tfrac{3\sqrt{3}}{4}\pi\cos\left(\tfrac{1}{4}\pi t\right)\\
\ddot{x}&=-\tfrac{3}{16}\pi^2\cos\left(\tfrac{1}{4}\pi t\right)-\tfrac{3\sqrt{3}}{16}\pi^2\sin\left(\tfrac{1}{4}\pi t\right)\\
&=-\tfrac{1}{16}\pi^2\left(3\cos\left(\tfrac{1}{4}\pi t\right)+3\sqrt{3}\sin\left(\tfrac{1}{4}\pi t\right)\right)\\
&=-\tfrac{1}{16}\pi^2 x
\end{align*}

This is of the form \(\ddot{x}=-\omega^2x\) with \(\omega=\frac{\pi}{4}\), so \(P\) is moving with simple harmonic motion about \(O\).

\item The period is
\[
T=\frac{2\pi}{\omega}=\frac{2\pi}{\frac{\pi}{4}}=8\,\mathrm{s}
\]

\item Write \(3\cos\left(\tfrac{1}{4}\pi t\right)+3\sqrt{3}\sin\left(\tfrac{1}{4}\pi t\right)\) as \(R\cos\left(\tfrac{1}{4}\pi t-\alpha\right)\), where \(R>0\) and \(0<\alpha<\frac{\pi}{2}\). Expanding and comparing coefficients:
\[
R\cos\alpha=3,\qquad R\sin\alpha=3\sqrt{3}
\]

so
\begin{align*}
R&=\sqrt{3^2+\left(3\sqrt{3}\right)^2}=\sqrt{36}=6\\
\tan\alpha&=\sqrt{3}\quad\Longrightarrow\quad\alpha=\frac{\pi}{3}
\end{align*}

Hence \(x=6\cos\left(\tfrac{1}{4}\pi t-\tfrac{1}{3}\pi\right)\), which oscillates between \(-6\) and \(6\), so the amplitude is \(6\,\mathrm{m}\).

\item For simple harmonic motion with amplitude \(a\), the maximum speed is \(a\omega\), so
\[
v_{\max}=6\times\frac{\pi}{4}=\frac{3\pi}{2}\,\mathrm{m\,s^{-1}}
\]

\item For simple harmonic motion with amplitude \(a\), the speed \(v\) at displacement \(x\) from the centre satisfies
\[
v^2=\omega^2\left(a^2-x^2\right)
\]

The speed is greater than \(\frac{3\sqrt{3}}{4}\pi\,\mathrm{m\,s^{-1}}\) when \(v^2>\left(\frac{3\sqrt{3}}{4}\pi\right)^2\). Substituting \(\omega=\frac{\pi}{4}\) and \(a=6\):
\begin{align*}
\frac{\pi^2}{16}\left(36-x^2\right)&>\left(\frac{3\sqrt{3}}{4}\pi\right)^2=\frac{27\pi^2}{16}\\
36-x^2&>27\\
x^2&<9
\end{align*}

so when \(-3<x<3\).

In each complete oscillation, \(P\) crosses the region \(-3<x<3\) twice, once in each direction, and by symmetry each crossing takes twice the time taken to travel from \(O\) to \(x=3\). The required time is therefore four times the time from \(O\) to \(x=3\).

Measuring time \(\tau\) seconds from an instant when \(P\) passes through \(O\) moving in the positive direction, the displacement is \(x=6\sin\left(\tfrac{1}{4}\pi\tau\right)\). Setting \(x=3\):
\begin{align*}
6\sin\left(\tfrac{1}{4}\pi\tau\right)&=3\\
\sin\left(\tfrac{1}{4}\pi\tau\right)&=\frac{1}{2}\\
\tfrac{1}{4}\pi\tau&=\frac{\pi}{6}\quad\Longrightarrow\quad\tau=\frac{2}{3}
\end{align*}

The total time is
\[
4\times\frac{2}{3}=\frac{8}{3}\,\mathrm{s}\approx 2.67\,\mathrm{s}
\]
\end{enumerate}
\end{djsSolution}

\newpage

% ---- Question 2 ----
\questionitem
% [Source: _QBT___Solns__Simple_Harmonic_Motion, Q2]

A particle $P$ of mass $0.5\,\mathrm{kg}$ moves in a straight line with simple harmonic motion about a centre $O$. The time taken for one complete oscillation is $\frac{\pi}{2}$ seconds. \\
When $P$ is $0.3\,\mathrm{m}$ from $O$, its speed is $1.6\,\mathrm{m\,s^{-1}}$. \\
\begin{enumerate}[label=\textbf{(\alph*)}]
\item Show that the amplitude of the motion is $0.5\,\mathrm{m}$. \marks{4} \\
\item Find the length of time within each complete oscillation for which the speed of $P$ is greater than $1\,\mathrm{m\,s^{-1}}$, giving your answer as an exact multiple of $\pi$. \marks{6} \\
\item Find the magnitude of the resultant force acting on $P$ at an instant when the kinetic energy of $P$ is $0.36\,\mathrm{J}$. \marks{4} \\
\end{enumerate}

\vspace*{10pt}

\begin{djsSolution}
\begin{enumerate}[label=\textbf{(\alph*)}, start=1]
\item Let the amplitude be \(a\,\mathrm{m}\) and the angular frequency \(\omega\,\mathrm{rad\,s^{-1}}\). The period of the motion is
\[
T=\dfrac{2\pi}{\omega}
\]

so
\[
\dfrac{2\pi}{\omega}=\dfrac{\pi}{2}\quad\Longrightarrow\quad \omega=4
\]

For simple harmonic motion with amplitude \(a\), the speed \(v\) at displacement \(x\) from the centre satisfies
\[
v^2=\omega^2\left(a^2-x^2\right)
\]

Substituting \(x=0.3\) and \(v=1.6\):
\begin{align*}
1.6^2&=4^2\left(a^2-0.3^2\right)\\
2.56&=16\left(a^2-0.09\right)\\
a^2-0.09&=0.16\\
a^2&=0.25\\
a&=0.5
\end{align*}

So the amplitude of the motion is \(0.5\,\mathrm{m}\), as required.

\item The speed is greater than \(1\,\mathrm{m\,s^{-1}}\) when \(v^2>1\), that is, when
\begin{align*}
16\left(0.25-x^2\right)&>1\\
0.25-x^2&>\dfrac{1}{16}\\
x^2&<\dfrac{3}{16}
\end{align*}

So the speed exceeds \(1\,\mathrm{m\,s^{-1}}\) exactly when \(P\) is within \(\frac{\sqrt{3}}{4}\,\mathrm{m}\) of \(O\).

Measure \(t\) from an instant when \(P\) passes through \(O\) moving in the positive direction. Then
\[
x=a\sin\omega t=0.5\sin 4t
\]

\(P\) first reaches \(x=\frac{\sqrt{3}}{4}\) when
\begin{align*}
0.5\sin 4t&=\dfrac{\sqrt{3}}{4}\\
\sin 4t&=\dfrac{\sqrt{3}}{2}\\
4t&=\dfrac{\pi}{3}\\
t&=\dfrac{\pi}{12}
\end{align*}

So \(P\) takes \(\frac{\pi}{12}\,\mathrm{s}\) to travel from \(O\) to a point \(\frac{\sqrt{3}}{4}\,\mathrm{m}\) from \(O\). The motion is symmetric, so in each of the four quarter-oscillations (from \(O\) out to an end point, and from an end point back to \(O\), on each side of \(O\)) \(P\) spends \(\frac{\pi}{12}\,\mathrm{s}\) within \(\frac{\sqrt{3}}{4}\,\mathrm{m}\) of \(O\). The total time within each complete oscillation is
\[
4\times\dfrac{\pi}{12}=\dfrac{\pi}{3}\,\mathrm{s}
\]

\item The kinetic energy is \(\frac12 mv^2\), so
\begin{align*}
\dfrac12\times 0.5\times v^2&=0.36\\
v^2&=1.44
\end{align*}

Using \(v^2=\omega^2\left(a^2-x^2\right)\):
\begin{align*}
1.44&=16\left(0.25-x^2\right)\\
0.25-x^2&=0.09\\
x^2&=0.16\\
x&=\pm 0.4
\end{align*}

So \(P\) is \(0.4\,\mathrm{m}\) from \(O\). In simple harmonic motion the acceleration is \(\ddot{x}=-\omega^2x\), so its magnitude here is
\[
\omega^2\times 0.4=16\times 0.4=6.4\,\mathrm{m\,s^{-2}}
\]

Using Newton's second law, the magnitude of the resultant force is
\[
F=0.5\times 6.4=3.2\,\mathrm{N}
\]
\end{enumerate}
\end{djsSolution}

\newpage

% ---- Question 3 ----
\questionitem
% [Source: _QBT___Solns__Simple_Harmonic_Motion, Q3]

A particle $P$ moves in a straight line with simple harmonic motion about a fixed point $O$. \\
When $P$ is $0.5\,\mathrm{m}$ from $O$, its speed is $3.6\,\mathrm{m\,s^{-1}}$, and when $P$ is $1.2\,\mathrm{m}$ from $O$, its speed is $1.5\,\mathrm{m\,s^{-1}}$. \\
The amplitude of the motion is $a$ metres. \\
\begin{enumerate}[label=\textbf{(\alph*)}]
\item Show that $a=1.3$. \marks{5} \\
\item Find the exact period of the motion and the greatest magnitude of the acceleration of $P$. \marks{3} \\
\end{enumerate}
During one complete oscillation, the magnitude of the acceleration of $P$ is at least $6.3\,\mathrm{m\,s^{-2}}$ for a total of $T$ seconds. \\
\begin{enumerate}[resume, label=\textbf{(\alph*)}]
\item Find the value of $T$. \marks{6} \\
\end{enumerate}

\vspace*{10pt}

\begin{djsSolution}
\begin{enumerate}[label=\textbf{(\alph*)}, start=1]
\item For simple harmonic motion with amplitude \(a\), the speed \(v\) at displacement \(x\) from the centre satisfies
\[
v^2=\omega^2\left(a^2-x^2\right)
\]

Substituting the two given positions:
\begin{align*}
3.6^2&=\omega^2\left(a^2-0.5^2\right) &&\Longrightarrow\quad 12.96=\omega^2\left(a^2-0.25\right)\\
1.5^2&=\omega^2\left(a^2-1.2^2\right) &&\Longrightarrow\quad 2.25=\omega^2\left(a^2-1.44\right)
\end{align*}

Subtracting the second equation from the first:
\begin{align*}
10.71&=\omega^2\left(1.44-0.25\right)\\
10.71&=1.19\,\omega^2\\
\omega^2&=9
\end{align*}

so \(\omega=3\). Substituting into the first equation:
\begin{align*}
12.96&=9\left(a^2-0.25\right)\\
a^2-0.25&=1.44\\
a^2&=1.69
\end{align*}

Since \(a>0\), \(a=1.3\), as required.

\item The period is
\[
\frac{2\pi}{\omega}=\frac{2\pi}{3}\,\mathrm{s}
\]

The acceleration is \(\ddot{x}=-\omega^2x\), so its magnitude is greatest at the extreme positions, where \(|x|=a\):
\[
\omega^2a=9\times1.3=11.7\,\mathrm{m\,s^{-2}}
\]

\item The magnitude of the acceleration is \(9|x|\), so it is at least \(6.3\,\mathrm{m\,s^{-2}}\) when
\[
9|x|\ge 6.3\quad\Longrightarrow\quad |x|\ge 0.7
\]

Measure time from an instant when \(P\) is at an extreme position on the positive side of \(O\), so that
\[
x=1.3\cos 3t
\]

Moving from this extreme position towards \(O\), \(P\) has \(x\ge 0.7\) until the time \(t_1\) given by
\begin{align*}
1.3\cos 3t_1&=0.7\\
\cos 3t_1&=\frac{7}{13}\\
t_1&=\frac{1}{3}\arccos\left(\frac{7}{13}\right)
\end{align*}

The motion is symmetrical about each extreme position and about \(O\), so in one complete oscillation \(P\) spends a time \(t_1\) with \(|x|\ge0.7\) on each side of each extreme position. This is four intervals of length \(t_1\):
\[
T=4t_1=\frac{4}{3}\arccos\left(\frac{7}{13}\right)\approx 1.34
\]

\end{enumerate}
\end{djsSolution}

\newpage

% ---- Question 4 ----
\questionitem
% [Source: _QBT___Solns__Simple_Harmonic_Motion, Q4]

The depth of water at the entrance to a harbour varies with the tides, and the water level moves with simple harmonic motion. \\
On a particular day, high tide occurs at 05:40, when the depth of water is $11.2\,\mathrm{m}$, and the next low tide occurs at 11:55, when the depth of water is $3.6\,\mathrm{m}$. \\
\begin{enumerate}[label=\textbf{(\alph*)}]
  \item Between 05:40 and 11:55 there are two moments at which the water level is falling at a rate of $1.5\,\mathrm{m\,h^{-1}}$. Find the depth of water at each of these moments. \marks{6} \\
\end{enumerate}
A ship can only use the harbour when the depth of water at the entrance is at least $5.5\,\mathrm{m}$. \\
\begin{enumerate}[label=\textbf{(\alph*)}, resume]
  \item Find the length of time between 05:40 and the following high tide during which the ship cannot use the harbour. \marks{4} \\
\end{enumerate}

\vspace*{10pt}

\begin{djsSolution}
The centre of the motion is the mean of the high and low tide depths, and the amplitude is half their difference:
\begin{align*}
\text{centre depth}&=\frac{11.2+3.6}{2}=7.4\,\mathrm{m}\\
a&=\frac{11.2-3.6}{2}=3.8\,\mathrm{m}
\end{align*}

From high tide at 05:40 to low tide at 11:55 is 6 hours 15 minutes, which is half a period, so the period is \(T=12.5\) hours. Using \(T=\frac{2\pi}{\omega}\):
\[
\omega=\frac{2\pi}{12.5}=\frac{4\pi}{25}\,\mathrm{rad\,h^{-1}}
\]

Let \(x\,\mathrm{m}\) be the height of the water level above its centre level of \(7.4\,\mathrm{m}\).
\begin{enumerate}[label=\textbf{(\alph*)}, start=1]
  \item For simple harmonic motion with amplitude \(a\), the speed \(v\) at displacement \(x\) from the centre satisfies
\[
v^2=\omega^2\left(a^2-x^2\right)
\]

Between high tide and low tide the water level is falling throughout, so the two moments are where the speed is \(1.5\,\mathrm{m\,h^{-1}}\). Substituting \(v=1.5\), \(a=3.8\) and \(\omega=\frac{4\pi}{25}\):
\begin{align*}
1.5^2&=\left(\frac{4\pi}{25}\right)^2\left(3.8^2-x^2\right)\\
x^2&=14.44-\frac{2.25\times625}{16\pi^2}\\
x^2&\approx 5.535\\
x&\approx \pm 2.353
\end{align*}

The depths of water are therefore
\[
7.4+2.353\approx 9.75\,\mathrm{m}\qquad\text{and}\qquad 7.4-2.353\approx 5.05\,\mathrm{m}
\]

  \item Let \(t\) be the time in hours after 05:40. At \(t=0\) the water is at high tide, so \(x=a\) when \(t=0\), and
\[
x=3.8\cos\left(\frac{4\pi t}{25}\right)
\]

The following high tide is at \(t=12.5\). The ship cannot use the harbour when the depth is less than \(5.5\,\mathrm{m}\), that is when \(x<5.5-7.4=-1.9\). The boundary values satisfy
\begin{align*}
3.8\cos\left(\frac{4\pi t}{25}\right)&=-1.9\\
\cos\left(\frac{4\pi t}{25}\right)&=-\frac{1}{2}
\end{align*}

For \(0\le t\le 12.5\), the angle \(\frac{4\pi t}{25}\) runs from \(0\) to \(2\pi\), so
\[
\frac{4\pi t}{25}=\frac{2\pi}{3}\quad\text{or}\quad\frac{4\pi t}{25}=\frac{4\pi}{3}\quad\Longrightarrow\quad t=\frac{25}{6}\quad\text{or}\quad t=\frac{25}{3}
\]

Between these times \(\cos\left(\frac{4\pi t}{25}\right)<-\frac12\), so the depth is below \(5.5\,\mathrm{m}\) (from 09:50 to 14:00). The length of time is
\[
\frac{25}{3}-\frac{25}{6}=\frac{25}{6}\text{ hours}
\]

This is 4 hours 10 minutes.
\end{enumerate}
\end{djsSolution}

\newpage

% ---- Question 5 ----
\questionitem
% [Source: _QBT___Solns__Simple_Harmonic_Motion, Q5]

A particle $P$ moves along the $x$-axis with simple harmonic motion. At time $t$ seconds, its displacement from the origin $O$ is $x$ metres, where
\[
\frac{\mathrm{d}^2x}{\mathrm{d}t^2}=9(2-x)
\]
and $P$ is instantaneously at rest when $x=5$. \\
\begin{enumerate}[label=\textbf{(\alph*)}]
\item Show that $v$, the velocity of $P$ in $\mathrm{m\,s^{-1}}$ at time $t$, satisfies
\[
v^2=9(5-x)(x+1)
\]
\marks[-20pt]{5}
\item When $t=0$, $P$ is at the point where $x=\frac{1}{2}$ and is moving in the positive $x$-direction. \\
Hence, or otherwise, show that
\[
x=2+3\sin\left(3t-\frac{\pi}{6}\right)
\]
\marks[-30pt]{6}
\end{enumerate}

\vspace*{10pt}

\begin{djsSolution}
\begin{enumerate}[label=\textbf{(\alph*)}, start=1]
\item Writing the acceleration as \(v\frac{\mathrm{d}v}{\mathrm{d}x}\):
\[
v\frac{\mathrm{d}v}{\mathrm{d}x}=9(2-x)=18-9x
\]

Separating the variables and integrating:
\begin{align*}
\int v\,\mathrm{d}v&=\int (18-9x)\,\mathrm{d}x\\
\tfrac12 v^2&=18x-\tfrac92x^2+c
\end{align*}

When \(x=5\), \(v=0\), so
\[
0=90-\tfrac{225}{2}+c\quad\Longrightarrow\quad c=\tfrac{45}{2}
\]

Hence
\begin{align*}
\tfrac12 v^2&=18x-\tfrac92x^2+\tfrac{45}{2}\\
v^2&=36x-9x^2+45\\
&=9\left(5+4x-x^2\right)\\
&=9(5-x)(x+1)
\end{align*}

as required.

\item From part (a), completing the square:
\[
(5-x)(x+1)=5+4x-x^2=9-(x-2)^2
\]

so \(v^2=9\left(9-(x-2)^2\right)\). At \(t=0\), \(P\) is moving in the positive \(x\)-direction, so \(v>0\) and
\[
\frac{\mathrm{d}x}{\mathrm{d}t}=3\sqrt{9-(x-2)^2}
\]

Separating the variables:
\[
\int\frac{1}{\sqrt{9-(x-2)^2}}\,\mathrm{d}x=\int 3\,\mathrm{d}t
\]

The formula booklet gives
\[
\int\frac{1}{\sqrt{a^2-x^2}}\,\mathrm{d}x=\arcsin\left(\frac{x}{a}\right)+K\qquad(|x|<a)
\]

Using this with \(a=3\) and \(x\) replaced by \(x-2\):
\[
\arcsin\left(\frac{x-2}{3}\right)=3t+K
\]

When \(t=0\), \(x=\frac12\), so
\[
K=\arcsin\left(\frac{\frac12-2}{3}\right)=\arcsin\left(-\tfrac12\right)=-\frac{\pi}{6}
\]

Therefore
\begin{align*}
\arcsin\left(\frac{x-2}{3}\right)&=3t-\frac{\pi}{6}\\
\frac{x-2}{3}&=\sin\left(3t-\frac{\pi}{6}\right)\\
x&=2+3\sin\left(3t-\frac{\pi}{6}\right)
\end{align*}

as required.
\end{enumerate}
\end{djsSolution}

\newpage

% ---- Question 6 ----
\questionitem
% [Source: _QBT___Solns__Simple_Harmonic_Motion, Q6]

A box of mass $m\,\mathrm{kg}$ rests on a horizontal platform. The platform moves vertically with simple harmonic motion of angular frequency $\omega\,\mathrm{rad\,s^{-1}}$ and amplitude $a\,\mathrm{m}$ about a mean position $O$. At time $t$ seconds, the upward displacement of the platform from $O$ is $x\,\mathrm{m}$. \\
The box remains in contact with the platform throughout the motion, and the greatest magnitude of the force exerted by the platform on the box is $\frac{7}{3}$ times the least magnitude of this force. \\
\begin{enumerate}[label=\textbf{(\alph*)}]
\item Show that $\omega^2=\frac{2g}{5a}$. \marks{3} \\
\end{enumerate}
The angular frequency is unchanged, but the amplitude of the motion is increased to $3a\,\mathrm{m}$. \\
\begin{enumerate}[label=\textbf{(\alph*)}, resume]
\item Show that the box leaves the platform, and find, in terms of $a$, the height above $O$ at which this first happens. \marks{3} \\
\item Find, in terms of $a$ and $g$, the speed of the box at the instant it leaves the platform. \marks{3} \\
\end{enumerate}

\vspace*{10pt}

\begin{djsSolution}
\begin{enumerate}[label=\textbf{(\alph*)}, start=1]
\item The box stays in contact with the platform, so it moves with the platform. For simple harmonic motion with angular frequency \(\omega\), the acceleration is
\[
\ddot{x}=-\omega^2x
\]

Let the force exerted by the platform on the box be \(R\,\mathrm{N}\). Take the direction of increasing \(x\) (upwards) as positive.

Using Newton's second law, with acceleration \(-\omega^2x\):
\begin{align*}
R-mg&=-m\omega^2x\\
R&=m\left(g-\omega^2x\right)
\end{align*}

Since \(-a\le x\le a\), \(R\) is greatest at the lowest point, \(x=-a\), and least at the highest point, \(x=a\):
\[
R_{\max}=m\left(g+\omega^2a\right),\qquad R_{\min}=m\left(g-\omega^2a\right)
\]

The greatest force is \(\frac{7}{3}\) times the least, so
\begin{align*}
3m\left(g+\omega^2a\right)&=7m\left(g-\omega^2a\right)\\
3g+3a\omega^2&=7g-7a\omega^2\\
10a\omega^2&=4g\\
\omega^2&=\frac{2g}{5a}
\end{align*}

as required.

\item While the box is in contact with the platform, the result of part (a) still holds:
\[
R=m\left(g-\omega^2x\right)=m\left(g-\frac{2gx}{5a}\right)
\]

The box leaves the platform when \(R=0\), which gives
\begin{align*}
g-\frac{2gx}{5a}&=0\\
x&=\frac{5a}{2}
\end{align*}

The platform now moves between \(x=-3a\) and \(x=3a\), and \(\frac{5a}{2}<3a\), so the platform reaches \(x=\frac{5a}{2}\) on its way up. For \(x>\frac{5a}{2}\) the formula would give \(R<0\), which is impossible, since the platform can only push on the box. Hence the box leaves the platform, and this first happens at a height of \(\frac{5a}{2}\,\mathrm{m}\) above \(O\).

\item Up to this instant the box moves with the platform, whose amplitude is \(3a\). For simple harmonic motion with amplitude \(A\), the speed \(v\) at displacement \(x\) from the centre satisfies
\[
v^2=\omega^2\left(A^2-x^2\right)
\]

Substituting \(\omega^2=\frac{2g}{5a}\), \(A=3a\) and \(x=\frac{5a}{2}\):
\begin{align*}
v^2&=\frac{2g}{5a}\left(9a^2-\frac{25a^2}{4}\right)\\
&=\frac{2g}{5a}\times\frac{11a^2}{4}\\
&=\frac{11ga}{10}\\
v&=\sqrt{\frac{11ga}{10}}
\end{align*}

The speed of the box as it leaves the platform is \(\sqrt{\frac{11ga}{10}}\,\mathrm{m\,s^{-1}}\).
\end{enumerate}
\end{djsSolution}

\newpage

% ---- Question 7 ----
\questionitem
% [Source: _QBT___Solns__Simple_Harmonic_Motion, Q7]

Throughout this question, use $g = 10\,\mathrm{m\,s^{-2}}$. \\
A light elastic string has natural length $0.9\,\mathrm{m}$ and modulus of elasticity $36\,\mathrm{N}$. \\
A particle $P$ of mass $0.4\,\mathrm{kg}$ is attached to one end of the string. The other end of the string is attached to a fixed point $A$. Particle $P$ hangs freely in equilibrium with $P$ vertically below $A$. \\
The particle is then pulled vertically down to a point $B$ and released from rest. \\
\begin{enumerate}[label=\textbf{(\alph*)}]
\item Show that, while the string is taut, $P$ moves with simple harmonic motion with period $\frac{\pi}{5}$ seconds. \marks{6} \\
\end{enumerate}
When the string first becomes slack, the speed of $P$ is $\sqrt{3}\,\mathrm{m\,s^{-1}}$. \\
\begin{enumerate}[resume, label=\textbf{(\alph*)}]
\item Find the distance $AB$. \marks{3} \\
\end{enumerate}
After the string becomes slack, $P$ continues to move vertically upwards until it first comes to instantaneous rest at the point $C$. \\
\begin{enumerate}[resume, label=\textbf{(\alph*)}]
\item Find the distance $BC$ and the exact time taken for $P$ to move from $B$ to $C$. \marks{6} \\
\end{enumerate}

\vspace*{10pt}

\begin{djsSolution}
\begin{enumerate}[label=\textbf{(\alph*)}, start=1]
\item Hooke's law is
\[
T=\frac{\lambda e}{l}
\]

where \(e\) is the extension and \(l\) the natural length. Here \(\lambda=36\) and \(l=0.9\), so \(T=40e\).

In equilibrium, let the extension be \(e_0\). Equating the tension to the weight gives
\[
40e_0=0.4\times 10=4\quad\Longrightarrow\quad e_0=0.1\,\mathrm{m}
\]

Let \(x\) be the displacement of \(P\) below the equilibrium position. While the string is taut, its extension is \(0.1+x\), so the tension is \(40(0.1+x)\).

Take the downward direction as positive. Using Newton's second law, with acceleration \(\ddot{x}\):
\begin{align*}
0.4\ddot{x}&=0.4\times 10-40(0.1+x)\\
0.4\ddot{x}&=4-4-40x\\
\ddot{x}&=-100x
\end{align*}

This is of the form \(\ddot{x}=-\omega^2x\) with \(\omega=10\), so while the string is taut \(P\) moves with simple harmonic motion about the equilibrium position. The period is
\[
\frac{2\pi}{\omega}=\frac{2\pi}{10}=\frac{\pi}{5}\,\mathrm{s}
\]

as required.

\item The string becomes slack when its extension is zero, which is when \(x=-0.1\). Let the amplitude be \(a\). Since \(P\) is released from rest at \(B\), \(B\) is at \(x=a\).

For simple harmonic motion with amplitude \(a\), the speed \(v\) at displacement \(x\) from the centre satisfies
\[
v^2=\omega^2\left(a^2-x^2\right)
\]

Substituting \(v=\sqrt{3}\), \(\omega=10\) and \(x=-0.1\):
\begin{align*}
3&=100\left(a^2-0.01\right)\\
a^2&=0.04\\
a&=0.2
\end{align*}

So \(B\) is \(0.2\,\mathrm{m}\) below the equilibrium position, which is \(0.9+0.1=1\,\mathrm{m}\) below \(A\):
\[
AB=0.9+0.1+0.2=1.2\,\mathrm{m}
\]

\item Once the string is slack, the only force on \(P\) is its weight, so \(P\) moves freely under gravity with constant acceleration. Take the upward direction as positive, so the acceleration is \(\alpha=-10\), the initial velocity \(u=\sqrt{3}\) and the final velocity \(v=0\).

The constant acceleration formula used is
\[
v^2=u^2+2\alpha s
\]

Substituting:
\[
0=3-20s\quad\Longrightarrow\quad s=0.15
\]

So \(C\) is \(0.15\,\mathrm{m}\) above the point where the string becomes slack. That point is \(0.2+0.1=0.3\,\mathrm{m}\) above \(B\), so
\[
BC=0.3+0.15=0.45\,\mathrm{m}
\]

For the time from \(B\) to the point where the string becomes slack, measure \(t\) from the release at \(B\). Since \(P\) starts from rest at \(x=a\), the motion is
\[
x=0.2\cos 10t
\]

The string becomes slack when \(x=-0.1\):
\begin{align*}
0.2\cos 10t&=-0.1\\
\cos 10t&=-\frac{1}{2}\\
10t&=\frac{2\pi}{3}\\
t&=\frac{\pi}{15}
\end{align*}

taking the first positive solution.

Let the time spent moving freely under gravity be \(t_2\). The constant acceleration formula used is
\[
v=u+\alpha t
\]

Substituting \(t=t_2\):
\[
0=\sqrt{3}-10t_2\quad\Longrightarrow\quad t_2=\frac{\sqrt{3}}{10}
\]

The total time from \(B\) to \(C\) is
\[
\frac{\pi}{15}+\frac{\sqrt{3}}{10}=\frac{2\pi+3\sqrt{3}}{30}\,\mathrm{s}
\]
\end{enumerate}
\end{djsSolution}

\newpage

% ---- Question 8 ----
\questionitem
% [Source: _QBT___Solns__Simple_Harmonic_Motion, Q8]

\begin{figure}[H]
    \centering
\begin{tikzpicture}[scale=1.3]
\pgfmathsetmacro{\ang}{asin(3/5)}
\def\planelen{7.2}
\def\groundlen{6.4}
\def\posA{0.42}
\def\posO{0.9}
\def\gap{0.07}
\def\prad{1.8pt}

\coordinate (base) at (0,0);
\coordinate (top) at (\ang:\planelen);
\coordinate (horiz) at (\groundlen,0);
\coordinate (A) at ($(base)!\posA!(top)$);
\coordinate (O) at ($(base)!\posO!(top)$);
\coordinate (P) at ($(A)+({\ang+90}:\gap)$);
\coordinate (Os) at ($(O)+({\ang+90}:\gap)$);

\draw[thick] (base) -- (horiz);
\draw[thick] (base) -- (top);

\draw[thick] (P) -- (Os) -- (O);

\fill (P) circle (\prad);

\pic["$\alpha$", draw, angle radius=10mm, angle eccentricity=1.35] {angle=horiz--base--top};

\node[above left] at (P) {$P$};
\node[below right] at (A) {$A$};
\node[below right] at (O) {$O$};
\end{tikzpicture}
\end{figure}


A particle $P$ of mass $m$ is attached to one end of a light elastic string of natural length $5a$ and modulus of elasticity $3mg$. The other end of the string is fixed to a point $O$ on a smooth plane, which is inclined at an angle $\alpha$ to the horizontal, where $\sin\alpha=\frac{3}{5}$. The string lies along a line of greatest slope of the plane and $P$ rests in equilibrium on the plane at a point $A$ below $O$, as shown in the diagram above. \\
At time $t=0$, $P$ is projected from $A$ with speed $2\sqrt{\frac{3ag}{5}}$ up the line of greatest slope towards $O$. At time $t$, the displacement of $P$ from $A$ is $x$, where the positive direction of $x$ is up the plane. \\
\begin{enumerate}[label=\textbf{(\alph*)}]
\item Show that, while the string is taut, $x$ satisfies the differential equation
\[
\frac{\mathrm{d}^2x}{\mathrm{d}t^2}+\frac{3gx}{5a}=0
\] \marks[-30pt]{6}
\item Find, in terms of $a$ and $g$, the time at which the string first becomes slack. \marks{5} \\
\item Find, in terms of $a$, the distance from $A$ of the highest point on the plane reached by $P$. \marks{3} \\
\end{enumerate}

\vspace*{10pt}

\begin{djsSolution}
\begin{enumerate}[label=\textbf{(\alph*)}, start=1]
\item Hooke's law is
\[
T=\frac{\lambda e}{l}
\]

where \(e\) is the extension and \(l\) the natural length. Here \(\lambda=3mg\) and \(l=5a\).

Let the extension at \(A\) be \(e_0\). Resolving parallel to the plane for the equilibrium at \(A\):
\begin{align*}
\frac{3mg e_0}{5a}&=mg\sin\alpha\\
\frac{3mg e_0}{5a}&=\frac{3mg}{5}\\
e_0&=a
\end{align*}

When the displacement of \(P\) from \(A\) is \(x\) up the plane, the extension is \(a-x\), so while the string is taut the tension, acting up the plane, is
\[
T=\frac{3mg(a-x)}{5a}
\]

The component of the weight down the plane is \(mg\sin\alpha=\frac{3mg}{5}\), and the normal reaction has no component along the plane.

Using Newton's second law, with acceleration \(\frac{\mathrm{d}^2x}{\mathrm{d}t^2}\):
\begin{align*}
m\frac{\mathrm{d}^2x}{\mathrm{d}t^2}&=\frac{3mg(a-x)}{5a}-\frac{3mg}{5}\\
m\frac{\mathrm{d}^2x}{\mathrm{d}t^2}&=\frac{3mg}{5}-\frac{3mgx}{5a}-\frac{3mg}{5}\\
\frac{\mathrm{d}^2x}{\mathrm{d}t^2}&=-\frac{3gx}{5a}\\
\frac{\mathrm{d}^2x}{\mathrm{d}t^2}+\frac{3gx}{5a}&=0
\end{align*}

as required.

\item By part (a), while the string is taut \(P\) moves with simple harmonic motion about \(A\), with
\[
\omega^2=\frac{3g}{5a}\quad\Longrightarrow\quad \omega=\sqrt{\frac{3g}{5a}}
\]

At \(t=0\), \(P\) is at the centre of the motion, so its speed \(2\sqrt{\frac{3ag}{5}}\) is the maximum speed. Let the amplitude be \(R\). The maximum speed is \(R\omega\), so
\begin{align*}
R\sqrt{\frac{3g}{5a}}&=2\sqrt{\frac{3ag}{5}}\\
R&=2\sqrt{\frac{3ag}{5}}\times\sqrt{\frac{5a}{3g}}\\
R&=2a
\end{align*}

Since \(x=0\) when \(t=0\) and \(P\) starts moving in the positive direction, the displacement is
\[
x=2a\sin(\omega t)
\]

The extension is \(a-x\), so the string first becomes slack when \(x=a\). Since \(a<2a\), \(P\) does reach this point:
\begin{align*}
2a\sin(\omega t)&=a\\
\sin(\omega t)&=\frac{1}{2}
\end{align*}

The first time this happens is when \(\omega t=\frac{\pi}{6}\), so
\[
t=\frac{\pi}{6}\sqrt{\frac{5a}{3g}}
\]

\item For simple harmonic motion with amplitude \(R\), the speed \(v\) at displacement \(x\) from the centre satisfies
\[
v^2=\omega^2\left(R^2-x^2\right)
\]

The string becomes slack when \(x=a\). Substituting \(\omega^2=\frac{3g}{5a}\), \(R=2a\) and \(x=a\):
\[
v^2=\frac{3g}{5a}\left(4a^2-a^2\right)=\frac{9ag}{5}
\]

After this the string is slack, and \(P\) moves a further distance \(d\) up the plane before coming to rest.

The work--energy principle is
\[
KE_0 + GPE_0 + EPE_0 + \sum W_{\text{other}} = KE_1 + GPE_1 + EPE_1
\]

where state 0 is the point where the string becomes slack and state 1 is the highest point. Take the initial position as the zero level for gravitational potential energy. The string is slack from state 0 to state 1, so \(EPE_0=EPE_1=0\). The normal reaction is perpendicular to the motion, so no other force does work, and \(\sum W_{\text{other}}=0\). This gives
\begin{align*}
\frac{1}{2}m\times\frac{9ag}{5}&=mgd\sin\alpha\\
\frac{9mag}{10}&=\frac{3mgd}{5}\\
d&=\frac{3a}{2}
\end{align*}

The distance from \(A\) of the highest point is
\[
a+\frac{3a}{2}=\frac{5a}{2}
\]
\end{enumerate}
\end{djsSolution}

\newpage

% ---- Question 9 ----
\questionitem
% [Source: _QBT___Solns__Simple_Harmonic_Motion, Q9]

\begin{figure}[H]
    \centering
\begin{tikzpicture}[scale=1]
\def\AB{10}
\def\AE{3}
\pgfmathsetmacro{\CE}{2/3}
\def\ext{0.6}
\def\tick{0.15}
\def\ctick{0.1}
\def\dimLow{0.7}
\def\dimHigh{1.5}
\def\prad{2.5pt}

\coordinate (A) at (0,0);
\coordinate (B) at (\AB,0);
\coordinate (E) at (\AE,0);
\coordinate (C) at ($(E)+(-\CE,0)$);

\draw[thick] ($(A)+(-\ext,0)$) -- ($(B)+(\ext,0)$);

\draw[thick] ($(A)+(0,-\tick)$) -- ($(A)+(0,\tick)$);
\draw[thick] ($(B)+(0,-\tick)$) -- ($(B)+(0,\tick)$);
\draw[thick] ($(C)+(0,-\ctick)$) -- ($(C)+(0,\ctick)$);

\fill[black] (E) circle (\prad);

\draw[{Stealth[length=2mm]}-{Stealth[length=2mm]}] ($(C)+(0,\dimLow)$) -- ($(E)+(0,\dimLow)$)
  node[midway, above] {$\frac{2}{3}L$};
\draw[{Stealth[length=2.5mm]}-{Stealth[length=2.5mm]}] ($(A)+(0,\dimHigh)$) -- ($(B)+(0,\dimHigh)$)
  node[midway, above] {$10L$};

\node[below=3pt] at (A) {$A$};
\node[below=3pt] at (B) {$B$};
\node[below=3pt] at (C) {$C$};
\node[below=3pt] at (E) {$E$};
\node[above right=3pt] at (E) {$P$};
\end{tikzpicture}
\end{figure}


A particle $P$ of mass $m$ lies on a smooth horizontal table. \\
One end of a light elastic string, of natural length $2L$ and modulus of elasticity $3mg$, is attached to $P$, and the other end is attached to a fixed point $A$ on the table. \\
One end of a second light elastic string, of natural length $4L$ and modulus of elasticity $2mg$, is also attached to $P$, and the other end is attached to a fixed point $B$ on the table, where $AB=10L$. \\
The particle $P$ rests in equilibrium at the point $E$, where $AEB$ is a straight line, as shown in the diagram above. \\
\begin{enumerate}[label=\textbf{(\alph*)}]
\item Show that $AE=3L$. \marks{2} \\
\end{enumerate}
The particle $P$ is now held at the point $C$ on $AE$, where $CE=\frac{2}{3}L$, and released from rest. \\
\begin{enumerate}[label=\textbf{(\alph*)}, resume]
\item Show that $P$ moves with simple harmonic motion with centre $E$ and period $\pi\sqrt{\frac{2L}{g}}$. \marks{6} \\
\item Find, in terms of $L$ and $g$, the speed of $P$ at the instant when $AP=\frac{5}{2}L$. \marks{3} \\
\item Find, in terms of $L$ and $g$, the exact length of time in each complete oscillation for which the tension in the string $AP$ is greater than $2mg$. \marks{6} \\
\end{enumerate}

\vspace*{10pt}

\begin{djsSolution}
\begin{enumerate}[label=\textbf{(\alph*)}, start=1]
\item Hooke's law is
\[
T=\frac{\lambda x}{l}
\]

where \(x\) is the extension and \(l\) the natural length.

At \(E\), the string \(AP\) has extension \(AE-2L\), and the string \(BP\) has length \(10L-AE\), so its extension is \(6L-AE\). The tensions are
\begin{align*}
T_{AP}&=\frac{3mg(AE-2L)}{2L}\\
T_{BP}&=\frac{2mg(6L-AE)}{4L}=\frac{mg(6L-AE)}{2L}
\end{align*}

In equilibrium the two tensions are equal, so
\begin{align*}
3(AE-2L)&=6L-AE\\
4AE&=12L\\
AE&=3L
\end{align*}

as required.

\item Let \(y\) be the displacement of \(P\) from \(E\), measured in the direction from \(A\) to \(B\). Then \(AP=3L+y\), with extension \(L+y\), and \(BP=7L-y\), with extension \(3L-y\). While both strings are taut, the tensions are
\begin{align*}
T_{AP}&=\frac{3mg(L+y)}{2L}\\
T_{BP}&=\frac{2mg(3L-y)}{4L}=\frac{mg(3L-y)}{2L}
\end{align*}

Take the direction from \(A\) to \(B\) as positive. The tension in \(BP\) acts towards \(B\) and the tension in \(AP\) acts towards \(A\).

Using Newton's second law, with acceleration \(\ddot{y}\):
\begin{align*}
m\ddot{y}&=\frac{mg(3L-y)}{2L}-\frac{3mg(L+y)}{2L}\\
&=\frac{mg}{2L}\left(3L-y-3L-3y\right)\\
&=-\frac{2mg}{L}\,y
\end{align*}

so
\[
\ddot{y}=-\frac{2g}{L}\,y
\]

This is of the form \(\ddot{y}=-\omega^2 y\) with \(\omega^2=\frac{2g}{L}\), so \(P\) moves with simple harmonic motion with centre \(E\), where \(y=0\).

\(P\) is released from rest at \(C\), where \(y=-\frac{2}{3}L\), so the amplitude is \(\frac{2}{3}L\) and \(-\frac{2}{3}L\le y\le\frac{2}{3}L\) throughout the motion. The extension of \(AP\) is then at least \(L-\frac{2}{3}L=\frac{1}{3}L>0\), and the extension of \(BP\) is at least \(3L-\frac{2}{3}L=\frac{7}{3}L>0\), so both strings stay taut and the equation of motion holds throughout.

The period is
\[
\frac{2\pi}{\omega}=2\pi\sqrt{\frac{L}{2g}}=\pi\sqrt{\frac{4L}{2g}}=\pi\sqrt{\frac{2L}{g}}
\]

as required.

\item When \(AP=\frac{5}{2}L\), the displacement is \(y=\frac{5}{2}L-3L=-\frac{1}{2}L\). The amplitude is \(a=\frac{2}{3}L\). For simple harmonic motion with amplitude \(a\), the speed \(v\) at displacement \(y\) from the centre satisfies
\[
v^2=\omega^2\left(a^2-y^2\right)
\]

Substituting \(\omega^2=\frac{2g}{L}\):
\begin{align*}
v^2&=\frac{2g}{L}\left(\frac{4L^2}{9}-\frac{L^2}{4}\right)\\
&=\frac{2g}{L}\times\frac{7L^2}{36}\\
&=\frac{7gL}{18}
\end{align*}

so the speed of \(P\) is
\[
v=\sqrt{\frac{7gL}{18}}=\frac{\sqrt{14gL}}{6}
\]

\item The tension in \(AP\) is greater than \(2mg\) when
\begin{align*}
\frac{3mg(L+y)}{2L}&>2mg\\
3(L+y)&>4L\\
y&>\frac{1}{3}L
\end{align*}

Measure the time \(t\) from the moment of release. \(P\) starts from rest at \(y=-\frac{2}{3}L\), so
\[
y=-\frac{2}{3}L\cos\omega t
\]

The condition \(y>\frac{1}{3}L\) becomes
\begin{align*}
-\frac{2}{3}L\cos\omega t&>\frac{1}{3}L\\
\cos\omega t&<-\frac{1}{2}
\end{align*}

During one complete oscillation, \(0\le\omega t\le 2\pi\), this holds for
\[
\frac{2\pi}{3}<\omega t<\frac{4\pi}{3}
\]

an interval of \(\frac{2\pi}{3}\) in \(\omega t\). With \(\omega=\sqrt{\frac{2g}{L}}\), the length of time is
\[
\frac{2\pi}{3}\times\sqrt{\frac{L}{2g}}=\frac{\pi}{3}\sqrt{\frac{4L}{2g}}=\frac{\pi}{3}\sqrt{\frac{2L}{g}}
\]
\end{enumerate}
\end{djsSolution}

\newpage

% ---- Question 10 ----
\questionitem
% [Source: _QBT___Solns__Simple_Harmonic_Motion, Q10]

A particle $P$ of mass $m$ is attached to one end of a light elastic string of natural length $l$ and modulus of elasticity $mg$. The other end of the string is attached to a fixed point $O$. \\
Particle $P$ is projected vertically downwards from $O$ with speed $\sqrt{gl}$. \\
\begin{enumerate}[label=\textbf{(\alph*)}]
\item Show that the greatest length of the string in the subsequent motion is $4l$. \marks{5} \\
\end{enumerate}
The string first becomes taut at time $t=0$. At time $t \ge 0$, the length of the string is $l+x$, where $x$ is the extension in the string. \\
\begin{enumerate}[label=\textbf{(\alph*)},resume]
\item Show that, while the string is taut, $x$ satisfies the differential equation
\[
\ddot{x}+\omega^2 x=g, \qquad \text{where } \omega^2=\dfrac{g}{l}
\]
\marks[-30pt]{3}
\item By using the substitution $x=y+\frac{g}{\omega^2}$, solve the differential equation and hence find, in terms of $l$ and $g$, the exact length of time for which the string is taut before it first becomes slack. \marks{12} \\
\end{enumerate}

\vspace*{10pt}

\begin{djsSolution}
\begin{enumerate}[label=\textbf{(\alph*)}, start=1]
\item Until the string becomes taut, \(P\) falls freely a distance \(l\) with acceleration \(g\). The constant acceleration formula used is
\[
v^2=u^2+2as
\]

Substituting \(u=\sqrt{gl}\), \(a=g\) and \(s=l\), where \(v\) is the speed when the string becomes taut:
\[
v^2=gl+2gl=3gl
\]

Let the greatest extension of the string be \(d\). At the lowest point \(P\) is instantaneously at rest. The work--energy principle is
\[
KE_0 + GPE_0 + EPE_0 + \sum W_{\text{other}} = KE_1 + GPE_1 + EPE_1
\]

where state 0 is when the string first becomes taut and state 1 is the lowest point. The elastic potential energy is
\[
EPE=\frac{kx^2}{2l}
\]

where \(k\) is the modulus of elasticity, \(x\) is the extension and \(l\) the natural length; here \(k=mg\). Take the final position as the zero level for gravitational potential energy. No other force does work, so \(\sum W_{\text{other}}=0\). Substituting:
\begin{align*}
\tfrac12 m(3gl)+mgd+0+0&=0+0+\frac{mgd^2}{2l}\\
3l^2+2ld&=d^2\\
d^2-2ld-3l^2&=0\\
(d-3l)(d+l)&=0
\end{align*}

Since \(d>0\), \(d=3l\), so the greatest length of the string is \(l+3l=4l\), as required.

\item Hooke's law is
\[
T=\frac{kx}{l}
\]

where \(k\) is the modulus of elasticity, \(x\) is the extension and \(l\) the natural length. Here \(k=mg\), so while the string is taut the tension is \(\frac{mgx}{l}\), acting upwards on \(P\), and the weight \(mg\) acts downwards. Take the direction of increasing \(x\) (downwards) as positive. Using Newton's second law, with acceleration \(\ddot{x}\):
\begin{align*}
m\ddot{x}&=mg-\frac{mgx}{l}\\
\ddot{x}+\frac{g}{l}x&=g
\end{align*}

so \(\ddot{x}+\omega^2x=g\), where \(\omega^2=\frac{g}{l}\), as required.

\item Since \(\frac{g}{\omega^2}=g\times\frac{l}{g}=l\), the substitution is \(x=y+l\), so \(\ddot{x}=\ddot{y}\). Substituting into the differential equation:
\begin{align*}
\ddot{y}+\omega^2(y+l)&=g\\
\ddot{y}+\omega^2y+g&=g\\
\ddot{y}+\omega^2y&=0
\end{align*}

The auxiliary equation is
\[
\lambda^2+\omega^2=0\quad\Longrightarrow\quad\lambda=\pm i\omega
\]

so
\begin{align*}
y&=A\cos\omega t+B\sin\omega t\\
x&=l+A\cos\omega t+B\sin\omega t
\end{align*}

Differentiating:
\[
\dot{x}=-A\omega\sin\omega t+B\omega\cos\omega t
\]

When \(t=0\), the string first becomes taut, so \(x=0\), and from part (a) the speed is \(\sqrt{3gl}\) downwards, so \(\dot{x}=\sqrt{3gl}\). Using \(x=0\) when \(t=0\):
\[
0=l+A\quad\Longrightarrow\quad A=-l
\]

Using \(\dot{x}=\sqrt{3gl}\) when \(t=0\):
\[
B\omega=\sqrt{3gl}\quad\Longrightarrow\quad B=\sqrt{3gl}\times\sqrt{\frac{l}{g}}=\sqrt{3}\,l
\]

Hence
\[
x=l-l\cos\omega t+\sqrt{3}\,l\sin\omega t
\]

Write \(\sqrt{3}\,l\sin\omega t-l\cos\omega t\) as \(R\sin(\omega t-\alpha)\), with \(R>0\) and \(0<\alpha<\frac{\pi}{2}\). Expanding:
\[
R\sin(\omega t-\alpha)=R\cos\alpha\sin\omega t-R\sin\alpha\cos\omega t
\]

Comparing coefficients:
\[
R\cos\alpha=\sqrt{3}\,l,\qquad R\sin\alpha=l
\]

so
\[
R=\sqrt{3l^2+l^2}=2l,\qquad \tan\alpha=\frac{1}{\sqrt{3}}\quad\Longrightarrow\quad\alpha=\frac{\pi}{6}
\]

Therefore the solution, valid while the string is taut, is
\[
x=l+2l\sin\left(\omega t-\frac{\pi}{6}\right),\qquad\text{where }\omega=\sqrt{\frac{g}{l}}
\]

The string becomes slack when \(x=0\):
\begin{align*}
l+2l\sin\left(\omega t-\frac{\pi}{6}\right)&=0\\
\sin\left(\omega t-\frac{\pi}{6}\right)&=-\frac12
\end{align*}

At \(t=0\), \(\omega t-\frac{\pi}{6}=-\frac{\pi}{6}\), and \(\omega t-\frac{\pi}{6}\) increases with \(t\). For \(-\frac{\pi}{6}<\omega t-\frac{\pi}{6}<\frac{7\pi}{6}\), \(\sin\left(\omega t-\frac{\pi}{6}\right)>-\frac12\), so \(x>0\) and the string stays taut. The next solution is therefore
\begin{align*}
\omega t-\frac{\pi}{6}&=\frac{7\pi}{6}\\
\omega t&=\frac{4\pi}{3}\\
t&=\frac{4\pi}{3}\sqrt{\frac{l}{g}}
\end{align*}
\end{enumerate}
\end{djsSolution}

\end{enumerate}
\end{document}
